\par\noindent\textbf{(1)} 在 \(x\ll1\) 的非磁性情形，电偶极项 \(a_1\approx\frac{2i}{3}(\varepsilon_r-1)/(\varepsilon_r+2)\,x^3\)，具体正负 \(i\) 与时谐约定有关。相应磁偶极 \(b_1=O(x^5)\)，因此普通弱磁介质首先由电偶极主导。
\par\noindent\textbf{(2)} Mie总散射截面 \(\sigma_{\rm sca}=(2\pi/k^2)\sum_{\ell=1}^\infty(2\ell+1)(|a_\ell|^2+|b_\ell|^2)\)。保留 \(\ell=1\) 电偶极，得 \(\sigma_{\rm sca}\approx(6\pi/k^2)|a_1|^2=(8\pi/3)k^4a^6|(\varepsilon_r-1)/(\varepsilon_r+2)|^2\)，为Rayleigh六次方律。
\par\noindent\textbf{(3)} 若通过高折射率和尺寸调节使电、磁偶极幅度可比，后向散射振幅在最低阶与 \(a_1-b_1\) 成正比，前向散射与 \(a_1+b_1\) 成正比。因此 \(a_1=b_1\) 能抑制后向散射，属于Kerker型干涉条件；纯非磁Rayleigh极限一般不满足。
\par\noindent\textbf{(4)} 被动无吸收散射体满足 \(\sigma_{\rm ext}=(2\pi/k^2)\sum_\ell(2\ell+1)\operatorname{Re}(a_\ell+b_\ell)\)，同时 \(\sigma_{\rm ext}=\sigma_{\rm sca}\ge0\)。因此不能通过设置 \(a_1=-b_1\) 就在保留非零总散射的同时让所有前向散射严格消失；必须考虑光学定理、相位与高阶项。
